\begin{lemma}[Простая нижняя оценка (лекции)]
При \(3\le m\le n\): \(R(m,n)\ge m+n-1\).
\end{lemma}
\begin{proof}
Граф \(K_{m-1,n-1}\) имеет \(m+n-2\) вершин. Клик размера 3 нет (две из трёх вершин
попали бы в одну долю), а \(m\ge3\) — значит, нет и \(K_m\). Независимое множество лежит
в одной доле (вершины разных долей смежны), поэтому \(\alpha=n-1<n\).
\end{proof}
Конструкция «\(m-1\) непересекающихся клик \(K_{n-1}\)» даёт заметно лучшее
\(R(m,n)>(m-1)(n-1)\) (см. рис.~\ref{fig:12-chvatal} в билете~\ref{ticket:12}), но всё
ещё лишь полиномиальную оценку. Эрдёш показал, что рост экспоненциальный.
\begin{graphillustration}{Простая нижняя оценка при $m=3$, $n=4$: в $K_{2,3}$ (5 вершин) нет треугольника, а $\alpha=3<4$. Значит, $R(3,4)\ge6$ (истинное значение 9).}{fig:09-bipartite}
\begin{tikzpicture}[font=\small]
\node[gi vertex] (a1) at (0.8,1.2) {};
\node[gi vertex] (a2) at (2.2,1.2) {};
\node[gi vertex] (b1) at (0,0) {};
\node[gi vertex] (b2) at (1.5,0) {};
\node[gi vertex] (b3) at (3,0) {};
\foreach \a in {a1,a2} \foreach \b in {b1,b2,b3} \draw[gi edge] (\a) -- (\b);
\node[gi group,fit=(a1)(a2),label=right:{$m-1=2$}] {};
\node[gi group,fit=(b1)(b2)(b3),label=right:{$n-1=3$}] {};
\end{tikzpicture}
\end{graphillustration}


\begin{theorem}[Эрдёш, 1947]
При \(n\ge2\): \(R(n,n)\ge2^{n/2}\).
\end{theorem}
\begin{proof}
\(R(2,2)=2=2^{2/2}\); пусть \(n\ge3\). Зафиксируем \(\nu<2^{n/2}\) помеченных вершин и
покажем, что есть граф на них без \(K_n\) и без независимого множества размера \(n\).
Пусть \(\mathcal G(\nu)=2^{\binom\nu2}\) — число всех графов на этих вершинах,
\(T(\nu,n)\) — число графов, содержащих \(K_n\), \(P(\nu,n)\) — содержащих независимое
множество размера \(n\). Если \(T+P<\mathcal G\), искомый граф существует.

\emph{\(T=P\)}: дополнение — биекция между этими классами.

\emph{Оценка \(T\)}: выберем \(n\) вершин будущей клики (\(\binom{\nu}{n}\) способов), рёбра
между ними проведём, остальные \(\binom\nu2-\binom n2\) пар — произвольно. Каждый граф с
\(K_n\) учтён хотя бы раз:
\[
T(\nu,n)\le\binom{\nu}{n}2^{\binom\nu2-\binom n2}
\le\frac{\nu^n}{n!}\,2^{-\frac{n(n-1)}2}\,\mathcal G(\nu)
<\frac{2^{n^2/2}}{n!}\,2^{-\frac{n^2}2+\frac n2}\,\mathcal G(\nu)
=\frac{(\sqrt2)^n}{n!}\,\mathcal G(\nu).
\]
\emph{\((\sqrt2)^n/n!<\tfrac12\) при \(n\ge3\)}: при \(n=3\) это \(2\sqrt2/6\approx0{,}47\), а при
переходе \(n\to n+1\) выражение умножается на \(\sqrt2/(n+1)<1\).

Итак, \(T<\mathcal G/2\), \(P=T<\mathcal G/2\), \(T+P<\mathcal G\). Значит, \(R(n,n)>\nu\) для каждого
целого \(\nu<2^{n/2}\), откуда \(R(n,n)\ge2^{n/2}\).
\end{proof}
\begin{tcolorbox}[colback=cbGrayBack,colframe=cbGray!55]
\textbf{Вероятностный смысл.} Выберем граф на \(\nu\) вершинах случайно (каждое ребро
независимо с вероятностью \(\tfrac12\)). Доказано: вероятность увидеть \(K_n\) или
независимое \(n\)-множество меньше 1, значит, хороший граф существует, хотя
явно его построить никто не умеет.
\end{tcolorbox}
\begin{graphillustration}{Подсчёт в доказательстве Эрдёша: фиксируем $n$ вершин и проводим все $\binom n2$ рёбер между ними (сплошные); остальные пары (серые) выбираем произвольно. Получаем не больше $\binom\nu n2^{\binom\nu2-\binom n2}$ графов с кликой $K_n$.}{fig:09-counting}
\begin{tikzpicture}[font=\small]
\foreach \i in {0,...,7} \node[gi small] (v\i) at (90+45*\i:1.4) {};
\foreach \i in {0,...,7} \foreach \j in {0,...,7} {\ifnum\i<\j \draw[black!18] (v\i) -- (v\j); \fi}
\draw[gi c1] (v0) -- (v1) -- (v2) -- (v0);
\draw[gi c1] (v0) -- (v7) -- (v1); \draw[gi c1] (v2) -- (v7);
\node[gi group,fit=(v0)(v1)(v2)(v7),label=above:{$n$ вершин}] {};
\node[gi note,anchor=west] at (2.2,0) {$\binom\nu n$ выборов клики,\\$2^{\binom\nu2-\binom n2}$ — остальное};
\end{tikzpicture}
\end{graphillustration}


\subsection*{Усиления}
\begin{cor}
При \(2\le m\le n\): \(R(m,n)\ge R(m,m)\ge2^{m/2}\), т.\,е. \(R(m,n)\ge2^{\min(m,n)/2}\).
\end{cor}
\begin{cor}[Общий критерий из доказательства]
Если \(\binom{\nu}{n}2^{1-\binom n2}<1\), то \(R(n,n)>\nu\).
\end{cor}
\begin{proof}
Из этого условия \(T+P\le2\binom\nu n2^{\binom\nu2-\binom n2}<\mathcal G(\nu)\).
\end{proof}
\begin{cor}[Точная форма оценки Эрдёша]
\[
R(n,n)\ge\bigl(1+o(1)\bigr)\frac{n}{e\sqrt2}\,2^{n/2}.
\]
\end{cor}
\begin{proof}
Используем \(\binom\nu n\le\frac{\nu^n}{n!}\le\left(\frac{e\nu}{n}\right)^n\) (так как \(n!\ge(n/e)^n\)):
\[
\binom{\nu}{n}2^{1-\binom n2}\le2\left(\frac{e\nu}{n}\right)^n2^{-\frac{n^2}{2}+\frac n2}
=2\left(\frac{e\sqrt2\,\nu}{n\,2^{n/2}}\right)^n .
\]
Для фиксированного \(\varepsilon>0\) и \(\nu=\bigl\lfloor(1-\varepsilon)\frac{n}{e\sqrt2}2^{n/2}\bigr\rfloor\)
правая часть не больше \(2(1-\varepsilon)^n\to0\), т.\,е. при больших \(n\) меньше 1, и
\(R(n,n)>\nu\) по предыдущему следствию.
\end{proof}
\textbf{Итог.} \(\sqrt2^{\,n}\lesssim R(n,n)\lesssim4^n\); показатели \(\sqrt2\) и \(4\) с 1947 года
улучшены лишь незначительно (верхняя — в 2023 году до \((4-\varepsilon)^n\)).
